{\large\textbf{E2 – Hidden pattern (10 pts)}}

You are given a flat semi-transparent foil with a micro-pattern printed on its
surface that is invisible to the naked eye. The pattern consists of a large
number of identical sinusoids with amplitude $A$, running horizontally with
spatial period $\Lambda$, and vertically shifted by distance $d$ relative to
each other, as schematically shown in Fig. 1. Under a microscope, one can see
that the printed pattern is composed of strictly horizontal line segments, each
vertically displaced from its neighbours by a constant pitch $s$, as shown in
Fig. 2.

\begin{center}\includegraphics[width=0.39\linewidth]{figures/EuPhO_2025_Q5_fig1.png}\end{center}

\begin{center}Figure 1: Pattern (not to scale)\end{center}

\begin{center}\includegraphics[width=0.48\linewidth]{figures/EuPhO_2025_Q5_fig2.png}\end{center}

\begin{center}Figure 2: Pattern as seen under microscope\end{center}

\textbf{Equipment (see also Fig. 3)}

A\quad Semi-transparent foil with a micro-pattern printed on its surface.

B\quad Laser diode with wavelength $\lambda = (654 \pm 5)\,\mathrm{nm}$. The
laser diode can be focused to the desired distance by rotating the end cap with
a lens inside.

\textbf{Warning:} Do not completely unscrew the end cap! Inside, there is an
oriented lens and a spring. No replacement laser will be given if damaged or
disassembled.

C\quad Two 90-degree L-shaped steel planks serving as stands for the foil and
the laser diode. The foil can be fixed to one of the planks using the provided
small clips. The laser diode can be mounted to the other plank with a larger
colored clip or with the provided rubber band.

D\quad A sheet of paper with a printed goniometer – a polar coordinate frame
with 1-mm radial steps and angular divisions in degrees.

\begin{center}\includegraphics[width=0.55\linewidth]{figures/EuPhO_2025_Q5_fig3.jpg}\end{center}

Figure 3: Components A, B, C, H, and J arranged for the experiment.

E\quad A screen: the large surface of the box containing the experimental
materials. Empty the box and place it on the desk with its large surface
vertical.

F\quad Ruler.

G\quad Measuring tape.

H\quad Adhesive tape attached to the ruler. Use pieces of the tape to fix the
printed goniometer to the screen or to secure components to the table. You can
ask for more tape if needed.

I\quad Millimeter graph paper.

J\quad An $80\,\mathrm{mm}$ paper measuring scale with diagonal reference lines
that allow you to measure fractions of the main scale divisions, accurate to
$\pm 0.1\,\mathrm{mm}$.

\textit{Hint:} In all of your measurements you are free to draw or put marks on
the screen.

\textit{Important:} Assume that the surface of the experimental desk is flat,
and the screen is strictly perpendicular to the desk.

\textbf{Tasks (10.0 pts)}

Determine as precisely as possible:

a\quad The sinusoid period $\Lambda$. (2 pts)

b\quad The vertical offset $d$ of the neighbouring sinusoids (2 pts)

c\quad The sinusoid amplitude $A$ (3 pts)

d\quad The step height $s$ (3 pts)

In all of the tasks you are expected to:

1.\quad sketch a setup and/or rationalize a method for measuring the
corresponding quantities;

2.\quad report your measurements and calculations in a tabular form;

3.\quad estimate the desired quantities and their uncertainties graphically,
whenever reasonable.
